\documentclass[10pt,a4paper]{amsart}
\usepackage[margin=19mm]{geometry}
\usepackage{amsmath,amssymb,mathtools}
\usepackage[scr=rsfs,cal=euler]{mathalfa}
\usepackage{xfrac}
\usepackage[unicode,hidelinks,bookmarks=false]{hyperref}
\newcommand{\ie}{i.e.\ }
\newcommand{\cf}{cf.\ }
\newcommand{\resp}{respectively\ }
\newcommand{\Subsection}{\subsection}
\providecommand{\nobreakdash}{\nobreak\hbox{-}}
\setlength{\emergencystretch}{2em}
\multlinegap0pt
\newcommand{\ds}{\displaystyle}
\usepackage{amssymb}
\usepackage{enumerate}
\usepackage{tikz}
\usepackage{caption}
\newtheorem{claim}{Claim}
\newtheorem{theorem}{Theorem}
\title{Dominion of some graphs}
\author{Julian Allagan, Benkam Bobga}
\date{Source: arXiv:2512.24115v1 (CC BY 4.0)}
\begin{document}
\maketitle

\noindent\emph{Source and licence.} Dominion of some graphs. Version arXiv:2512.24115v1 (CC BY 4.0), distributed by arXiv under the Creative Commons Attribution 4.0 International licence, http://creativecommons.org/licenses/by/4.0/, as recorded in the arXiv licence metadata for this exact version and shown on the abstract page https://arxiv.org/abs/2512.24115v1. Full licence notice: \texttt{licenses/arxiv-2512.24115-CC-BY-4.0.txt}.\par\smallskip

\noindent\textit{Editorial scope.} Counted unit: the article's theorem T1 (source0.tex lines 171--179). The article's complete original proof of it is delivered with it (source0.tex lines 181--268, one \texttt{proof} environment of 88 lines). No mathematics is written, repaired or re-derived: the statement and the proof are the article's own lines, in the article's own notation, and no retained line is substituted or re-flowed.  Retained uncounted as the source-local context the proof needs: lines 162--164.  Nothing was omitted from any retained span, so the package carries no omission record.  No retained line contains a citation, so the wrapper carries no bibliography and no bibliography entry.  No retained line contains a figure environment, an image-inclusion command or drawing code, so no image is delivered and the package declares no asset dependency.  Editorial formatting only, in addition: an AMS \texttt{amsart} preamble that restates the article's own declarations and macros used by the retained lines, namely the environments \texttt{claim}, \texttt{theorem} on separate counters, together with the standard packages those lines need.  The retained environments print in this document as Claim 1, Theorem 1 in order of appearance, whereas the article prints the same items with the numbering of its own full document.  The complete native source of the article as distributed in the arXiv e-print for this exact version is bundled unchanged as \texttt{sources/191853/source0.tex}.
\begin{claim}{\label{c4}}
	When $n=3k, k\ge 1$, $\zeta(P_n)=1$ and the unique $\gamma$-set in $P_n$ contains $v_{n-1}$. Moreover, when $n=3k+2, k\ge 1$, there is only one $\gamma$-set in $P_n$ containing $v_{n}$. 
\end{claim}

\begin{theorem}{\label{T1}}
	If $G$ is a path on $n$ vertices, then its dominion is \\
	
	$\zeta(G)=\begin{cases} 
	1, & \text{if} \ n \equiv 0\pmod3 \hspace{.5in} (i)\\
	\frac{(n+2)(n+11)}{18}-1, &  \text{if} \ n \equiv 1\pmod3 \hspace{.5in} (ii)\\
	{\lceil \frac{n}{3} \rceil}+1, & \text{if} \ n \equiv 2\pmod3. \hspace{.5in} (iii)
	\end{cases}$
\end{theorem}

\begin{proof}
	Part (i). This follows directly from Claim \ref{c4}. %Let $G=P_{3k}$. Because $\gamma(G)=k$ and each vertex $v\in V(G)$ covers at most two other vertices, any $\gamma$-set, say $S$ must include $v_{3k-1}$. It follows from Claim \ref{c4} that $S$ is unique and $S=\{v_2,v_5,\ldots,v_{3k-1}\}$, which indicates that $\zeta(G)=1$, when $n \equiv 0\pmod3$. 
	\newline
	
	Part (iii).
   	Suppose that $n=3k+2$, $k\ge 1$, and $G=P_n$. We have $\gamma(G)=k+1$ and every $\gamma$-set $S\subseteq V(G)$ contains exactly one of $v_n$, $v_{n-1}$. From Claim \ref{c4} we know that exactly one such $S$ contains $v_n$, so our task is to show that exactly ${\lceil \frac{n}{3} \rceil}=k+1$ such $S$ contain $v_{n-1}$.
	Suppose that $S$ is a $\gamma$-set for $G$ containing $v_{n-1}$. By the reasoning in the proof of Claim \ref{c4}, each $3$-set $\ds T_r=\{v_{3r+1},v_{3r+2},v_{3r+3}\}$, $r=0,\ldots,k-1$, contains exactly one vertex of $S$. 
	\indent If $0\le r< k-1$ and $\ds S\cap T_r=\{v_{3r+1}\}$ then, because $S$ is dominating in $G$, $S\cap T_{r+1}=\{v_{3r+4}\}$; consequently, if $S\cap T_r=\{v_{3r+1}\}$ then $S\cap T_s=\{v_{3s+1}\}$, $r\le s \le k-1$. 
	
	\indent On the other hand, if $S\cap T_r=\{v_{3r+2}\}$ then $S\cap T_{r+1}$ is either $\{v_{3r+4}\}$ or $\{v_{3r+5}\}$. Since $S\cap T_{0}$ is either $\{v_1\}$ or $\{v_2\}$ (since $v_3$ does not cover $v_1$), it follows that the $\gamma$-set for $G=P_n$, $n=3k+2$, containing $v_{n-1}$ are the following:
	\begin{align}
	    S_0&=\{v_1,v_4,\ldots,v_{n-4},v_{n-1}\},\nonumber \\
	     S_1&=\{v_2,v_4,\ldots,v_{n-4},v_{n-1}\},\nonumber \\
	     \vdots \nonumber \\
	     S_r&=\{v_2,v_5,\ldots,v_{3r-1},v_{3r+1},\ldots,  v_{n-4},v_{n-1}\},\nonumber \\
	     \vdots \nonumber \\
	     S_k&=\{v_2,v_5,\ldots,v_{n-3},v_{n-1}\}, \nonumber
	\end{align}
and there are exactly $k+1$ of these sets.

	Part(ii)
	We recursively obtain the $\zeta$ values of $G_k=P_{3k+1}:=v_1-v_2-\ldots-v_{3k}-v_{3k+1}$ as we increase the value of $k\ge 1$. For simplicity, let $\zeta(G_k)=\zeta_k$, $k\ge 1$.
	
	When $k=1$, it is easy to verify that the $\gamma$-sets of $G_1$ are $\{v_1,v_3\}$, $\{v_1,v_4\}$,$\{v_2,v_3\}$, $\{v_2,v_4\}$. Hence $\zeta_1=4$.
	
	When $k=2$, it is clear that $\gamma(G_2)=3$. Since $v_6$ covers $v_5$ and $v_7$, we add $v_6$ to each of the $\gamma$-sets of $G_1$ to get $\{v_1,v_3, v_6\}$, $\{v_1,v_4, v_6\}$, $\{v_2,v_3, v_6\}$, $\{v_2,v_4, v_6\}$. Further, we add $v_7$ to each $\gamma$-set of $G_1$ that includes $v_4$ since $v_4$ covers $v_5$, and $v_7$ covers $v_6$. This produces two additional $\gamma$-sets for $G_2$, namely $\{v_1,v_4, v_7\}, \{v_2,v_4, v_7\}$. Finally, with the new and unique set $\{v_2, v_5\}$ which covers all other vertices except $v_7$, we generate two more $\gamma$-sets of $G_2$ using $v_6$ for one and $v_7$ for the other. It is clear that $v_6$ covers $v_7$ (and vice-versa), giving  $\{v_2,v_5, v_6\}, \{v_2, v_5, v_7\}$. Hence $\zeta_2=8$.
	
	When $k=3$, it is also clear that $\gamma(G_3)=4$. We use a very similar construction as in the case when $k=2$. Since $v_9$ covers $v_8$ and $v_{10}$, we add $v_9$ to each of the $8$ $\gamma$-sets of $G_2$ to get 
	$\{v_1,v_3, v_6, v_9\}$, $\{v_1,v_4, v_6, v_9\}$, $\{v_2,v_3, v_6, v_9\}$, $\{v_2,v_4, v_6, v_9\}$, $\{v_1,v_4, v_7, v_9\}, \{v_2,v_4, v_7, v_9\}$, $\{v_2,v_5, v_6, v_9\}, \{v_2, v_5, v_7, v_9\}$
	. Further, we add $v_{10}$ to each $\gamma$-set of $G_2$ that includes $v_7$ since $v_7$ covers $v_8$, and $v_{10}$ covers $v_{9}$. This produces three additional $\gamma$-sets for $G_3$, namely $\{v_1,v_4, v_7, v_{10}\}, \{v_2,v_4, v_7, v_{10}\}, \{v_2, v_5, v_7, v_{10}\}$. Finally, with the unique set $\{v_2, v_5, v_8\}$ which covers all other vertices except $v_{10}$, we generate two more $\gamma$-sets of $G_3$ using $v_9$ for one and $v_{10}$ for the other. Thus, $v_9$ covers $v_{10}$ (and vice-versa) giving  $\{v_2,v_5, v_8, v_9\}, \{v_2, v_5, v_8, v_{10}\}$. Hence $\zeta_3=13$. Now, we generalize the previous constructions for all $k\ge 2$.
	
	It is obvious that $\gamma(G_k)=k+1$. So, for $k\ge 2$, we repeatedly form the $\gamma$-sets of $G_k$ by extending those of $G_{k-1}$ as follows:
	
	\begin{enumerate}
		\item[(a)] Since $v_{3k}$ covers both $v_{3k-1}$ and $v_{3k+1}$, we add $v_{3k}$ to each $\gamma$-set of $G_{k-1}$.
		\item[(b)] We add $v_{3k+1}$ to each $\gamma$-set of $G_{k-1}$ that includes $v_{3k-2}$ since $v_{3k-2}$ covers $v_{3k-1}$, and $v_{3k+1}$ covers $v_{3k}$.
		\item[(c)] Given the unique set $\{v_2,v_5,\ldots,v_{3k-1}\}$ which covers all other vertices of $G_k$ except $v_{3k+1}$, we generate two more $\gamma$-sets of $G_k$ using $v_{3k}$ for one and $v_{3k+1}$ for the other. Clearly $v_{3k}$ covers $v_{3k+1}$ (and vice-versa) and we have $\{v_2,v_5,\ldots,v_{3k-1}, v_{3k}\}$ and $\{v_2,v_5,\ldots,v_{3k-1}, v_{3k+1}\}$. 
	\end{enumerate}
	From steps (a)-(c), we generate a sequence of $\zeta_k$ values, with $k\ge 1$, giving 
	$$4, 8, 13, 19, 26, 34, 43, \ldots  $$
	
	For $k\ge 1$ let $A_k=\{S\subseteq V(G_k)\ |\ S \ \text{is a}\ \gamma-\text{set in}\ G_k \ \text{and} \ v_{3k+1} \in S\}$,
	 $B_k=\{S\subseteq V(G_k)\ |\ S \ \text{is a}\ \gamma-\text{set in}\ G_k \ \text{and} \ v_{3k+1} \notin S\}$, and $a_k=|A_k|, b_k=|B_k|$. Then $\zeta(G_k)=a_k+b_k$. Notice that $S\in B_k \implies v_{3k}\in S$. If $S\in B_k$, there are $3$ possibilities, if $k\ge 2$:
	 \begin{enumerate}
	     \item $ v_{3k-1}, v_{3k-2} \notin S,$ in which case $S\setminus \{v_{3k}\}\in B_{k-1}.$
	      \item $ v_{3k-1} \in S.$ 
	      \item $ v_{3k-2} \in S.$ 
	 \end{enumerate}
	 
	\indent It is not possible that $\{  v_{3k-2}, v_{3k-1},  v_{3k}\subseteq S \in B_k\}$, because this inclusion would imply that $S\setminus \{v_{3k-1},  v_{3k}\} $ is a dominating set in $G_{k-1}$, which is impossible, since $|S\setminus \{v_{3k-1},  v_{3k}\}|=k-1<\gamma(G_{k-1})=k$. Therefore, for $S\in B_k$, $2$ and $3$ are disjoint possibilities.
	
	\indent In case $2$ we see that $S'=S\setminus \{v_{3k-1},  v_{3k}\}$ is a dominating set in $P_{3(k-1)}$, with $k-1=\gamma(P_{3(k-1)})$ vertices, and we know that there is only one such set, $S'=\{v_j\ | \ j\equiv 2\pmod 3, 2\le j \le 3k-4\}$.
	
	\indent In case $3$, $S\setminus \{  v_{3k}\}\in A_{k-1}$. Furthermore, if $S' \in A_{k-1}$ then $S'\cup\{v_{3k}\} \in B_k$. Similar converses apply to cases $1$ and $2$: in case $2$, if $S'$ is the unique $\gamma$-set in $P_{3(k-1)}$ then $S'\cup \{v_{3k-1},  v_{3k}\} \in B_k$, and in case 1, if $S' \in B_{k-1}$ then $S\cup \{v_{3k}\} \in B_k$. Putting these observations together, we have that $b_k=a_{k-1}+b_{k-1}+1 \ \ \ (*)$.% $\tag{$\star$}$ %{\label{eq}}$
	
	\indent If $S\in A_k$ then $|S\setminus \{v_{3k+1}\}|=k$; since each set $T_r$, $r=0, \ldots, k-1$ must contain at least one vertex of $S\setminus \{v_{3k+1}\}$, it follows that each $T_r$ contains exactly one vertex of $S$. It follows that $v_{3r+3}\notin S$, $r=o\ldots k-1$, for if $v_{3r+3}\in S$ then, because there cannot be more than $2$ vertices between consecutive elements of $S$, it would follow that $v_{3s+3}\in S$, $0\le s \le r$. But then $v_3$ is the sole representative of $S$ in $T_0$, which contradicts the premise that $S$ is a $\gamma$-set in $G_k$. 
	
	\indent Thus, $S\in A_k$ implies $v_{3k}\notin S$, so the sole vertex in $S\cap T_{k-1}$ is either $v_{3k-1}$ or $v_{3k-2}=v_{3(k-1)+1}$. If $v_{3k-1} \in S$, then $S\setminus \{v_{3k+1}\}$ is a $\gamma$-set in $P_{3k},$ so $S=\{v_t\ | \ t=2\pmod 3, \ 2\le t\le 3k-1\} \cup \{v_{3k+1}\}$. If $v_{3k-2} \in S$, then $S\setminus \{v_{3k+1}\} \in A_{k-1}$; conversely, if $S' \in A_{k-1}$ then $v_{3k-2} \in S'$ and $S'\cup \{v_{3k-2}\}\in A_k$.
	
	\indent From the preceding paragraph we conclude that, for $k\ge 2$, $a_k=1+a_{k-1}$. Since $a_1=2$, it follows that $a_k=k+1$, $k=1,2, \ldots$
	
	\indent From $(*)$, we have that, for $k\ge 2$,
	\begin{align}
	    b_k&=a_{k-1}+b_{k-1}+1 =b_{k-1}+k+1 \nonumber \\
	    &= b_1+[(k+1)+k+\ldots + 3] \nonumber \\
	    &=2+\frac{(k+2)(k+1)}2-3 \nonumber \\
	    &=\frac{k^2+3k}2. \nonumber
	\end{align} 
	Therefore, for $k\ge 2$,
	\begin{align}
	    \zeta(G_k)&=a_k+b_k=k+1+\frac{k^2+3k}2 \nonumber \\
	    &=\frac{k^2+5k+2}2 = \frac{(n+2)(n+11)}{18}-1 \nonumber
	\end{align}
	if $n=3k+1$.
	
	
	
	
	
	
% 	It is not too hard to see that the sequence follows the recursion:
% 	\begin{equation}\label{eq1}
% 	\zeta_{3k+1}=\zeta_{3(k-1)+1}+\zeta_{3(k-1)+2}+ \zeta_{3(k-1)+3}, \tag{$\star$}
% 	\end{equation} 
% 	where $\zeta_{3(k-1)+1}=\zeta(G_{k-1})$, the dominion of $G_{k-1}$, and $\zeta_{3(k-1)+3}=\zeta(P_{3k})=1$, $\zeta_{3(k-1)+2}=\zeta(P_{3(k-1)+2})={\lceil \frac{3(k-1)+2}{3} \rceil}+1$, according to Parts (i), (iii), respectively. % Using a numerical approximation, we obtain a general form of the sequence whose first seventh terms are shown 
% 	From the recursion (\ref{eq1})  we obtain that $\ds \zeta_k=\frac{1}2k^2+\frac{5}2k+1=\frac{(k+1)(k+4)}2-1$. Now, since $\ds k=\frac{n-1}3$, the result follows from a direct substitution.
	
\end{proof}

\end{document}
